\clearpage
\properlane{clothed-charity}{The Invited People Clothed in Charity}

God's saving promise forms a people whose life must answer the love of
the Son. The feast is not fulfilled simply by assembling a crowd.
The guests are called into communion with the Bridegroom, and charity
is the life appropriate to that communion. Augustine and Gregory
therefore make the garment a question about belonging: how can a
person stand among Christ's people and yet resist the love that gives
their common life its form?

\subsection{What possession cannot prove}

Augustine's argument begins with gifts Christians rightly value.
Baptism, coming to church, receiving at the altar, fasting, knowledge,
and extraordinary spiritual gifts all belong to the world of serious
religious life. He does not make them worthless. He asks whether
possessing any one of them is sufficient to distinguish the guest
whose presence bears fruit from the guest whose presence ends in
judgment. Since good and bad can share these outwardly, none alone
answers the garment's question. The distinction must reach the way
the gift is received and lived.\footnote{Augustine, Sermon 90.5--6.}

This reasoning prevents two contrary mistakes. Mere participation
cannot serve as a guarantee of final blessedness. Equally, the
failure of a participant does not make God's gift false. Augustine
asks why great goods may profit a person nothing, not whether those
goods have ceased to be good. His argument about charity belongs
inside sacramental life. It is addressed to those who have entered
the feast and must learn to honour it.

Even almsgiving and martyr-like endurance can become occasions of
self-display. Augustine follows Paul's account of charity to the
intention that animates an action. A gift to the poor can be performed
for glory; an impressive religious achievement can nourish the
desire to be admired. The garment cannot be reduced to a visible
list of actions any more than it can be reduced to a list of
possessions. Yet it must take form in actions. The heart's love and
the neighbour's actual good belong together.

The Introit places this examination under a promise that precedes
achievement. The Lord calls himself the people's salvation before
their response is heard. The appended psalm summons them to listen,
and its fuller history shows gifts received with different results.
Augustine's exposition of Psalm 77 includes faithful hearers among
the ancient people; grace is not the exclusive boast of those who
now hear the psalm in church. A Christian who turns Israel's failures
into proof of his own security has missed the force of the
instruction. He too is a recipient whose heart must answer.
\footnote{Augustine, \work{Enarrationes} 77.1--3.}

The Collect then supplies the proper response to such dependence.
Instead of announcing that every obstruction has been conquered,
the people ask God to remove what hinders their service. Body and
mind are to be made ready together. Charity is not purchased by
producing a flawless religious appearance before approaching God.
It grows within a life delivered, ordered, and sustained by him.
The prayer acknowledges both the reality of service and the need
for divine assistance in performing it.

\subsection{The Bridegroom gives the garment its meaning}

Gregory arrives at charity from another direction. The garment is
nuptial because the Lord himself came in love to unite the Church
to himself. Its meaning rests in the Bridegroom, not in an arbitrary
rule of admission. To come to the wedding without charity is to
enter the celebration of Christ's loving union while refusing the
very love displayed there.\footnote{Gregory, \work{Gospel Homilies}
38.9; compare 38.3.}

His explanation of the wedding is carefully Christological. He
speaks of the Incarnation and then of Christ's union with the
Church. The Incarnation is not a marriage of two persons, one
divine and one human. Gregory expressly rejects that conclusion:
Christ is one Person in two natures. The bridal image unfolds the
gift of the Son to his people; it cannot overturn the identity of
the Son who gives himself. Thus the allegory deepens the narrative
without making its human arrangements a literal account of the
divine mystery.

Gregory's weaving analogy gives this love its twofold demand.
As a garment is woven upon upper and lower pieces of wood, charity
requires love of God and love of neighbour. A supposed love of God
that abandons a neighbour's need is incomplete in its very form;
service of a neighbour that lets love of God go cold also fails
to answer the whole command. His further image of twice-dyed
scarlet presses the same relation: charity burns toward God and
toward the person beside us.\footnote{Gregory, \work{Gospel Homilies}
38.10. The analogy concerns the two supports of the weaving;
Gregory does not identify two threads.}

The Gradual gives this charity a form of prayer. In Augustine's
reading, the Church offers her prayer within the evening sacrifice
of her Head. Her love does not rise independently of Christ's
Passion. The Bridegroom's gift grounds her own offering, so that
prayer and charity cannot be rival ways of obtaining access to God.
The Church prays because she belongs to Christ; she must love in
the manner of the one in whom she prays.

The Alleluia carries that love outward. Bellarmine's explanation
of proclaiming God's deeds begins with the desire that others
should know the beloved. Augustine hears the evangelists in the
announcement among nations. The gathered people are not to treat
their knowledge as a possession that becomes less valuable when
shared. Praise seeks new hearers. The Gospel's servants go out
with a corresponding generosity: the feast is to be filled.
\footnote{Augustine, \work{Enarrationes} 140.2--4 and 104.1--3;
Bellarmine, \work{Psalms} 104:1--2.}

\subsection{The difficult neighbour tests the love}

Affection for those who already love us does not exhaust Augustine's
garment. He widens the examination to strangers and enemies.
Shared human origin makes another person a neighbour; the work
of Christ gathers those whom sin has scattered. He then asks
whether a Christian can desire the evil person's healing rather
than destruction. Christ and Stephen supply the pattern of prayer
for persecutors. The enemy's evil is opposed precisely by desiring
that he cease to be evil.\footnote{Augustine, Sermon 90.7--10.}

Gregory likewise says that charity is tested by adverse hatred.
Love of a friend must be in God, and love of an enemy for God's
sake. The question is not whether the enemy has earned warm
feeling. It concerns whether love remains directed toward his
good when he offers no answering goodwill. This reaches further
than courtesy toward congenial people and further than peaceful
avoidance of those one dislikes.\footnote{Gregory, \work{Gospel
Homilies} 38.11.}

The Epistle locates that demanding love within a common body.
Truthful speech respects the neighbour who depends on truthful
information. Anger placed under restraint refuses to turn a
present injury into the organizing principle of a relationship.
Honest labour becomes a gift to need. These commands give charity
visible form without making the visible form sufficient by itself.
A person can ask both whether he helped and whether he sought
the neighbour's good, rather than his own display.

The Offertory keeps this practice from resting on an illusion
that love makes hostility disappear. The singer still walks in
tribulation. Augustine locates the life secured by God beyond
what an enemy can take away. Such trust makes enemy-love possible
without calling injury unreal. The worshipper can refuse revenge
because his final safety does not depend on mastering the enemy
through hatred. God's saving hand and the Christian's generous
hands are related through trust, not through a promise that all
danger will end immediately.

This also gives patience within the mixed assembly a positive
content. Gregory knows the difficulty of living among people
whose actions cause grief. Patience is not approval of their
wrong. Charity refuses to let another person's wrong become the
reason for one's own surrender to hatred. The Gospel's king
remains judge; the neighbour remains someone whose conversion
can be desired.\footnote{Gregory, \work{Gospel Homilies} 38.7,
11; Augustine, \work{Enarrationes} 137.10--13.}

\subsection{Imperfect charity and the medicine of the feast}

Gregory's consolation is as necessary as his warning. Someone
who has the garment imperfectly is not in the same condition as
someone wholly without it. Real charity can be weak and require
growth. The king's coming need not produce despair in the one
who genuinely loves yet recognizes his own incompleteness.
Gregory interrupts the severe exposition to address precisely
that person before returning to the guest who lacks charity.
\footnote{Gregory, \work{Gospel Homilies} 38.12.}

Augustine also describes charity as something to be born,
nourished, and increased. Disordered desire remains a struggle;
the answer is the growth of charity until its perfection.
Consequently the distinction between a mixed feast now and a
perfect feast at the end is not permission to postpone conversion.
The present is the place in which growth occurs. It is also not
a demand that every weakness be read as proof of rejection.
The garment's presence has to be distinguished from its final
perfection.\footnote{Augustine, Sermon 90.2--3, 6.}

The Secret's \textit{nobis} places the congregation's actual
condition within the prayer over the gifts. Schuster's account
of sacramental fruit answers Augustine's question about gifts
shared by good and bad: the sacrament is real, and recipients
still need the dispositions through which its saving power
bears fruit. The Church does not solve the garment's question
by declaring sacred participation unnecessary. She brings the
question into her petition for saving participation.

The Communion's wish then gives a voice to imperfect charity.
God has commanded; the worshipper asks for directed ways.
Augustine hears prayer for enabling grace in this wish, while
Schuster hears delight in God's service. The movement can
include both a real love of obedience and a real need to be
helped. One need not pretend to have arrived in order to desire
the right destination.

The Postcommunion gathers these needs under medicine. A
worshipper who has heard Augustine ask for the enemy's healing
now asks to be healed himself. The words expose the ease with
which one can diagnose another's lack of charity while excusing
one's own perversity. The final prayer is in the plural. The
congregation seeks freedom and adherence together, as a people
still dependent on the Lord who called himself their salvation
at the beginning.\footnote{Schuster, \work{Sacramentary} III,
pp.~173--174; Augustine, \work{Enarrationes} 118.4--5; Sermon
90.9--10.}

\subsection{The king's judgment and the final feast}

The wideness of the invitation does not erase the king's
inspection. Augustine reads the single rejected guest as
representing many, and he also speaks of the good as many in
themselves though few in comparison. His purpose is to urge
worthy participation that reaches the feast from which evil
is excluded. Turning the scene into a numerical calculation
would miss his distinction between the two feasts and the
present summons to charity.\footnote{Augustine, Sermon 90.4.}

Gregory makes uncertainty about one's perseverance a reason
for humility. Good beginnings can be abandoned; a life long
misdirected can be turned through repentance. Today's
appearance cannot settle the final judgment. Hope rests in
divine mercy and calls for continued response. The prayer's
request for lasting adherence therefore has an eternal
seriousness: it asks that the charity appropriate to the
Son's feast endure to its fulfilment.

\begin{fourSenses}
\item[Literal.] A prepared royal feast meets refusal, further invitation,
mixed attendance, and judgment. The Epistle commands conduct fitting
members of one another; the orations ask for freedom, saving benefit,
and healing. The psalms give hearing, prayer, praise, confidence, and
desire their places in this response.
\item[Allegorical.] Augustine and Gregory recognize Christ's union with
the Church in the wedding. Charity honours the Bridegroom who came in
love. The Church's prayer is offered in her crucified Head, whose gift
forms the people he gathers.
\item[Moral.] Membership calls for truthful, generous, patient love,
including desire for an enemy's healing. Imperfect charity needs growth;
neither sacramental pride nor despair answers the invitation. The
worshipper asks to receive fruitfully and to live what he receives.
\item[Anagogical.] The present mixed assembly looks toward the feast
where charity is perfected and evil no longer shares the table. Judgment
summons perseverance; God's enduring lordship gives hope beyond temporary
safety. The final word sought by the people is lasting adherence to him.
\end{fourSenses}
